% Начинаем с новой страницы, заголовок центрирован, прописными буквами
\chapter*{РЕФЕРАТ}
\addcontentsline{toc}{chapter}{РЕФЕРАТ} % Добавляем в оглавление

% 1. Заголовочная часть (Метаданные)
\noindent Пояснительная записка: 75 с., 15 рис., 8 табл., 35 источников, 4 прил.

\vspace{0.5cm}

% Впишите свою тему ПРОПИСНЫМИ буквами
\noindent \MakeUppercase{РАЗРАБОТКА АВТОМАТИЗИРОВАННОЙ СИСТЕМЫ МОНИТОРИНГА ПАРАМЕТРОВ ОКРУЖАЮЩЕЙ СРЕДЫ}

\vspace{0.8cm}

% 2. Реферативная часть (Текст)
% Текст должен содержать: объект исследования, цель, методы, результаты и выводы.

Объект исследования --- система беспроводного сбора данных с датчиков температуры и влажности в режиме реального времени.

Целью дипломного проекта является проектирование и программная реализация программно-аппаратного комплекса для мониторинга климатических условий в серверных помещениях.

В процессе работы проводился анализ существующих решений на базе протокола MQTT. Выполнено проектирование принципиальной электрической схемы устройства на базе микроконтроллера ESP32, разработаны алгоритмы энергосбережения и веб-интерфейс для визуализации данных.

В результате выполнения проекта была создана рабочая модель устройства, обеспечивающая точность измерений $\pm 0.5$ градусов Цельсия и автономную работу до 12 месяцев. Проведено тестирование системы в условиях промышленного предприятия, подтвердившее стабильность передачи данных.

Область применения: системы автоматизации зданий («умный дом»), контроль микроклимата в медицинских и складских помещениях.

\vspace{1cm}

% 3. Ключевые слова (необязательно по СТП, но часто требуют кафедры)
\noindent Ключевые слова: МИКРОКОНТРОЛЛЕР, ДАТЧИК, МОНИТОРИНГ, БЕСПРОВОДНАЯ СЕТЬ, ИНТЕРФЕЙС, АЛГОРИТМ.

\newpage